% !TeX root = main.tex
\section{Preliminaries} \label{sec:prelim}
In this section we provide the definition of the model and problem, as well as a selection of helpful results and observations in the low-bandwidth model from prior works \cite{gupta-hirvonen-etal-2022-sparse-matrix-multiplication, gupta-korhonen-etal-2025-low-bandwidth-matrix, gupta-suomela-vahidi-2026-rectangular-matrix}.

\subsection{Model and Problem Definitions} \label{subsec:definitions}

We work in the \emph{low-bandwidth model}, a distributed system that consists of $n$ machines, denoted by
$M = \{1,2,\dots,n\}$, where all pairs of machines are directly connected.
Computation proceeds in synchronous communication rounds. And, for the purposes of this work, we only consider deterministic algorithms.

\subparagraph{Local computation.}
Each machine has unbounded local memory and local computation is free. The complexity measure of interest is the number of communication rounds.

\subparagraph{Communication.}
In each round, every machine may
\begin{itemize}
	\item send at most one $O(\log n)$-bit message to another machine, and
	\item receive at most one $O(\log n)$-bit message from another machine.
\end{itemize}
Equivalently, in each round the communication forms a directed graph of maximum in-degree and out-degree~$1$. An algorithm is invalid if $2$ messages are sent to a machine in any round.

\subparagraph{Problem.}
We study uniformly sparse matrix multiplication. Given two $n\times n$ matrices $A,B$, the task is to compute specified entries of their product $X=AB$.

The instance includes binary masks $\bar{A},\bar{B},\bar{X}\in\{0,1\}^{n\times n}$ that mark either allowed input entries or requested output entries; all other input entries are zero.
We may occasionally abuse notation and identify a mask with the set of its marked positions, and write
\[
  Q=\{(i,k)\in[n]\times[n]:\bar{X}_{ik}=1\}.
\]
A matrix or a mask is \emph{uniformly $d$-sparse} if each row and each column contains at most $d$ marked positions.
We denote the class of such masks by $\US(d)$ and assume that $\bar{A},\bar{B},\bar{X}\in\US(d)$, where $d\in[n]$.
Note that the full product $X$ need not be sparse; however, $\bar{X}$ is guaranteed to be.

Throughout the paper, matrix multiplication is performed over an arbitrary semiring $(\mathcal{S},+,\cdot)$. We assume that the input entries, products, and partial sums used by the algorithms admit encodings of size $O(\log n)$ bits.
In particular, each such value fits within one communication message in the low-bandwidth model.

\subparagraph{Input and output distribution.}
The input is evenly distributed among the $n$ machines, and each machine is responsible for producing an evenly sized portion of the output.
More concretely:
\begin{itemize}
  \item each machine initially holds at most $O(d)$ input entries in total;
  \item each machine is required to output at most $O(d)$ entries of $X$ requested by $\bar{X}$.
\end{itemize}
The algorithm may choose the balanced assignment of requested output positions to machines. 

\subparagraph{Supported setting.}
In the supported setting, the input masks, the output mask, and the assignment of input values to machines are fixed in advance and known to all machines.
The output assignment is chosen from this public information before the input values are revealed and is also known to all machines.
The numerical values at marked input positions are initially supplied to their holders and are not part of this public information.
All machines know $n$ and $d$.

\subparagraph{Complexity.}
An algorithm runs in $T(n,d)$ rounds if, after $T(n,d)$ communication
rounds, all machines have produced their assigned output.

\subsection{Observations and Toolbox}\label{subsec:toolbox}

For any integer $n>0$, let $[n]$ denote the set $\{1,\ldots,n\}$.


We next briefly mention some useful tools. Note that machines do not coordinate during the algorithm to determine which machine they need to send a message to, or which machine they need to receive one from. Instead, the scheduling is already hard-coded into the algorithm; during execution, the values are revealed and transferred. This avoids the costly runtime coordination overhead. 

\begin{lemma}[Broadcast~{\cite[Lemma 6.9]{gupta-korhonen-etal-2025-low-bandwidth-matrix}}]
	\label{lem:broadcast}
	In the low-bandwidth model, broadcasting a single bit $v \in \{0,1\}$ to all machines takes $\Theta(\log n)$ rounds. Moreover, if $k>1$ messages are given, broadcasting can be done in $O(k+ \log n)$ rounds.
\end{lemma}
\begin{proof}
	We can broadcast in a power-jumping manner. Suppose machine~$1$ initially holds the value $v$. Before round~$i$, the first $2^{i-1}$ machines know $v$. Each of them sends $v$ to a distinct machine that does not yet know it, until all machines have received it. For example, machine~$j$ sends to machine~$j+2^{i-1}$ whenever $j+2^{i-1}\leq n$. After round~$i$, $\min\{2^i,n\}$ machines know $v$, so all machines know it after $O(\log n)$ rounds.
	
	This argument works for any value $v$ that fits in a single message. By carefully interleaving broadcasts of this form~\cite{kannan-ballard-park-2016-a-high-performance-parallel}, we can schedule the messages in a pipeline and achieve $O(k+\log n)$ rounds.
	
	We refer the reader to \cite[Lemma 6.9]{gupta-korhonen-etal-2025-low-bandwidth-matrix} for details of the lower bound.
\end{proof}

\begin{lemma}[Aggregation~{\cite[Lemma 6.4]{gupta-korhonen-etal-2025-low-bandwidth-matrix}}]
	\label{lem:aggregation}
	In the low-bandwidth model where all $n$ machines are given a single bit in $\{0,1\}$, computing logical AND, OR, and sum of the values takes $\Theta(\log n)$ rounds. Moreover, if machines are instead given $k$ values that fit in a message, aggregation can be done in $O(k + \log n)$ rounds.
\end{lemma}
\begin{proof}
	The idea is the same as in \Cref{lem:broadcast}, with the communication done in reverse order. This also works for values that fit in a message, provided that the intermediate results fit in a message as well. The same pipelining idea of \cite{kannan-ballard-park-2016-a-high-performance-parallel} extends the argument to $k$ values.
	
	We refer the reader to \cite[Lemma 6.4]{gupta-korhonen-etal-2025-low-bandwidth-matrix} for details of the lower bound.
\end{proof}

We occasionally use the term \emph{distribution} $D$ of data to specify which machine $m\in M$ holds which value $v$, such as an element of a matrix. Typically, $v$ is small enough to fit in a single message, i.e., $O(\log n)$ bits.

\begin{lemma}[Redistribution \cite{gupta-suomela-vahidi-2026-rectangular-matrix}]\label{lem:redistribution-offline}
	Given globally known distributions $D$ and $D'$ of $d$ values per machine, we can rearrange the values from $D$ to $D'$ in $d$ rounds of the low-bandwidth model.
\end{lemma}
\begin{proof}
	Consider the bipartite multigraph of senders and receivers, where each value forms an edge between its home machine in $D$ and its destination machine in $D'$. This multigraph is $d$-regular and therefore admits a proper edge coloring with $d$ colors. We iterate over the color classes, sending all messages of one color class in each round. Since each color class is a matching, every machine sends and receives at most one message per round. As $D$ and $D'$ are globally known, all machines can compute the same coloring locally using a fixed deterministic procedure, without communication.
\end{proof}



